% Generated deterministically from the audited Markdown source.
% Responsible author: Carptopus.
\documentclass[11pt]{article}
\usepackage[a4paper,margin=27mm]{geometry}
\usepackage[T1]{fontenc}
\usepackage[utf8]{inputenc}
\usepackage{lmodern}
\usepackage{microtype}
\usepackage{amsmath,amssymb,amsthm,mathtools}
\usepackage{xcolor}
\usepackage[colorlinks=true,linkcolor=blue!55!black,citecolor=blue!55!black,urlcolor=blue!65!black]{hyperref}
\usepackage{xurl}
\usepackage{fancyhdr}

\newtheorem{theorem}{Theorem}
\numberwithin{equation}{section}
\pagestyle{fancy}
\fancyhf{}
\fancyhead[L]{Carptopus}
\fancyhead[R]{Connected Peck counterexamples}
\fancyfoot[C]{\thepage}
\setlength{\headheight}{14pt}
\setlength{\parindent}{0pt}
\setlength{\parskip}{0.42em}
\setlength{\emergencystretch}{4em}

\title{Connected Peck posets with\\non-log-concave antichain polynomials}
\author{Carptopus\\\href{mailto:carptopus@163.com}{carptopus@163.com}}
\date{Manuscript, 2 September 2026}

\hypersetup{
  pdftitle={Connected Peck posets with non-log-concave antichain polynomials},
  pdfauthor={Carptopus},
  pdfsubject={An infinite counterexample family to Ding and Dong Conjecture C},
  pdfkeywords={Peck poset, antichain polynomial, log-concavity, Sperner property, independence polynomial, counterexample}
}

\begin{document}
\maketitle

\begin{center}
\fcolorbox{black!35}{black!4}{\parbox{0.88\textwidth}{\centering\small
Internally verified candidate manuscript, version 0.1-beta. The mathematical claims have passed
project-internal zero-trust proof and manuscript audits; external review is pending.}}
\end{center}

\begin{abstract}


Ding and Dong conjectured that the antichain generating polynomial of every finite connected Peck
poset is log-concave. We give an infinite family of height-one counterexamples. For integers
$a,c\geq 1$, the poset $P_{a,c}$ has two equally sized ranks and a connected bipartite Hasse graph
with a perfect matching. We prove that it is Peck and that its antichain polynomial is

\[
\mathcal N_{P_{a,c}}(x)=(1+2x)^a+2(1+x)^{a+c}-2(1+x)^a.
\]

For every $a\geq 7$, the specialization $P_{a,2}$ violates log-concavity. The first member on this
parameter line has 18 elements and log-concavity defect $-72$.

\end{abstract}

\section{Construction and statement}


Ding and Dong state as Conjecture C that the antichain generating polynomial of every finite
connected Peck poset is log-concave. Here ``height one'' means that maximal chains have one cover
edge, or equivalently two rank levels; Ding and Dong number those levels as ranks 1 and 2.

Our construction was motivated by Bhattacharyya and Kahn's nonunimodal independence-polynomial
construction for a bipartite graph. Their two bipartition classes have unequal sizes, so that graph
does not directly give a rank-symmetric Peck poset. The balanced rank layers, the added perfect
matching, and the closed formula below are the features used here to obtain a connected Peck
counterexample family.

Let $a,c\geq 1$. Define a two-rank poset $P_{a,c}$ with lower and upper ranks

\[
L=A\sqcup W,\qquad R=B\sqcup C,
\]

where

\[
|A|=|C|=a,\qquad |W|=|B|=c.
\]

The cover relations are all pairs in $A\times B$, the matching pairs $A_i<C_i$, and all pairs in
$W\times(B\sqcup C)$.

\begin{theorem}


For all $a,c\geq 1$, the poset $P_{a,c}$ is a finite connected Peck poset and

\[
\mathcal N_{P_{a,c}}(x)=(1+2x)^a+2(1+x)^{a+c}-2(1+x)^a.
\tag{1}
\]

For every $a\geq 7$, the coefficient sequence of $\mathcal N_{P_{a,2}}(x)$ is not log-concave.
Consequently, the connected-Peck log-concavity conjecture is false.

\end{theorem}

\section{The Peck property}


The Hasse graph is connected because $a,c\geq 1$. Both ranks have size $a+c$, so the poset is
rank-symmetric and rank-unimodal. The edges

\[
\{A_iC_i:1\leq i\leq a\}\cup\{W_jB_j:1\leq j\leq c\}
\]

form a perfect matching. By K\H{o}nig's theorem, the maximum independent-set size of the bipartite
Hasse graph is $a+c$. Antichains are exactly independent sets of this graph, while either rank is
an antichain of size $a+c$; hence $P_{a,c}$ is Sperner.

Every element has a neighbour in the other rank, and no chain contains three elements. Thus all
maximal chains have two elements. For $k\geq 2$, the whole $2(a+c)$-element poset is a maximum
$k$-family and its size equals the sum of the two largest rank sizes. Therefore $P_{a,c}$ is
strongly Sperner. It is consequently Peck under the definition used by Ding and Dong.

\section{The antichain polynomial}


We count independent sets according to whether they meet $W$. If no vertex of $W$ is chosen, the
remaining graph consists of the complete bipartite graph between $A$ and $B$, together with the
matching from $A$ to $C$. Its independence polynomial is

\[
(1+2x)^a+(1+x)^{a+c}-(1+x)^a.
\]

Indeed, the first term counts choices that use no vertex of $B$, while the second counts choices
that use no vertex of $A$; their overlap consists of arbitrary subsets of $C$ and is exactly
$(1+x)^a$, which explains the subtraction.

If a nonempty subset of $W$ is chosen, no vertex of $B\sqcup C$ may be chosen, while the vertices
of $A$ remain arbitrary. This contributes

\[
((1+x)^c-1)(1+x)^a=(1+x)^{a+c}-(1+x)^a.
\]

Their sum is (1).

\section{An infinite family of failures}


Set $c=2$, and write $p_i=[x^i]\mathcal N_{P_{a,2}}(x)$. The top three coefficients are

\[
p_{a+2}=2,\qquad p_{a+1}=2(a+2),\qquad
p_a=2^a-2+(a+2)(a+1).
\]

Log-concavity at position $a+1$ would require

\[
2^{a+1}\leq 2(a+2)(a+3)+4.
\tag{2}
\]

Define

\[
F(a)=2^{a+1}-2(a+2)(a+3)-4.
\]

Then $F(7)=72>0$, and

\[
F(a+1)-F(a)=2^{a+1}-4(a+3)>0
\]

for every $a\geq 7$. Thus (2) fails for all $a\geq 7$.

At $a=7$, the polynomial is

\[
1+18x+114x^2+378x^3+742x^4+882x^5+602x^6+198x^7+18x^8+2x^9,
\]

and the final internal log-concavity defect is

\[
18^2-198\cdot2=-72.
\]

Direct enumeration of all $2^{18}$ vertex subsets agrees coefficientwise with (1). The same
verification for $a=1,\ldots,6$ shows that $a=7$ is the first failure on the $c=2$ parameter
line. We do not claim global minimality among all connected Peck posets.

\begin{thebibliography}{9}
\footnotesize
\raggedright

\bibitem{bhattacharyya-kahn}
A.~Bhattacharyya and J.~Kahn,
\emph{A bipartite graph with non-unimodal independent set sequence},
Electronic Journal of Combinatorics \textbf{20} (2013), Paper 11.
\href{https://www.combinatorics.org/ojs/index.php/eljc/article/view/v20i4p11}{journal page}.

\bibitem{ding-dong}
C.~Ding and P.~Dong,
\emph{Antichain generating polynomials of posets},
arXiv:1905.06692 (2019).
\href{https://arxiv.org/abs/1905.06692}{arXiv:1905.06692}.

\end{thebibliography}

\paragraph{AI-assisted research and writing disclosure.}
Carptopus is the author responsible for the mathematical claims and the public version of this
manuscript. OpenAI Codex was used as an AI-assisted tool for literature organization, proof
development and stress testing, exact verification, adversarial auditing, and manuscript
preparation.

\end{document}
